\documentclass[12pt]{article}
\usepackage[utf8]{inputenc}
\usepackage[bahasa]{babel}
\usepackage{xcolor}
\usepackage{soul}
\usepackage{xurl}


\sethlcolor{yellow}
\emergencystretch=10em

\begin{document}

\begin{flushleft}
\textbf{Nama: Riziq Lili Ulil Albab} \\
\textbf{NIM: 058487469}
\end{flushleft}

\vspace{0.5cm}


\section*{Jawaban}
\begin{enumerate}

\item Berdasarkan artikel di atas, jawablah pertanyaan berikut!

\begin{enumerate}

\item Tuliskan judul artikel yang tepat untuk tulisan/naskah tersebut!

\textbf{“Ancaman Punahnya Bahasa Daerah dan Nilai Luhur Peribahasa”}


\vspace{0.5cm}

\item Perbaikilah penulisan kutipan yang ditulis secara tidak tepat!

\begin{quote}
Memiliki banyak kekayaan merupakan impian setiap orang. Berdasarkan sumber keberadaannya kekayaan (materi) berasal dari dua sumber yaitu berasal dari: (1) warisan/peninggalan/pemberian dan (2) kekayaan yang dicari dengan segala daya. Indonesia sebagai sebuah negara kepulauan memiliki keduanya.

Kekayaan negara yang berupa warisan adalah kekayaan alam berupa tambang, kayu, air, flora dan fauna, beserta jumlah penduduknya. Kekayaan tersebut diberikan Sang Maha Pencipta yang Pengasih dan penyayang kepada setiap makhluk-Nya.

Indonesia memiliki ribuan pulau dengan jumlah suku bangsa sekitar 1340. \hl{Menurut Liputan6.com, jumlah tersebut merupakan jumlah suku bangsa terbesar kedua di dunia setelah India.} Dari 1340 suku bangsa tersebut terdapat 718 bahasa lokal/daerah. Besarnya jumlah bahasa daerah ini juga menduduki nomor dua terbesar di dunia setelah Papua Nugini. \hl{Nabila Oudri menyatakan bahwa dari 718 bahasa daerah di Indonesia, sebagian besar kondisinya terancam punah dan kritis.} \hl{Berdasarkan data dari Kementerian Pendidikan, Kebudayaan, Riset, dan Teknologi, sebanyak 25 bahasa daerah terancam punah, enam bahasa dalam kondisi kritis, dan 11 bahasa telah dinyatakan punah.}

Data di atas merupakan gambaran yang seharusnya membuat bangsa ini prihatin terhadap kondisi kekayaan bahasa daerah kita. Berikutnya yang perlu kita waspadai yaitu akan terkikisnya kekayaan intelektual Indonesia yang lahir bersama bahasa dan tata krama masyarakat setempat. Apa itu? Satu dari sekian banyak kekayaan intelektual yakni peribahasa, termasuk di dalamnya pepatah, bidal, pemeo, tamsil, dan kalimat-kalimat perumpamaan lain.

\hl{Menurut KBBI, peribahasa adalah ungkapan atau kalimat ringkas padat yang berisi perbandingan, perumpamaan, nasihat, prinsip hidup, atau aturan tingkah laku.} \hl{Yudhistira, Lenin, dan Irawan (2022:44--45) menyatakan bahwa peribahasa merupakan rumusan kebijaksanaan masyarakat yang menunjukkan sifat waspada dan \textit{eling} yang berkaitan dengan moral dan kebajikan hidup.} \hl{Selain itu, mereka juga menyatakan bahwa peribahasa dapat menjadi sarana untuk mengungkapkan isi batin dan sebagai penanda nilai luhur budaya.}

Dari sumber tersebut diketahui bahwa di dalam peribahasa terkandung nasihat, teguran, atau sindiran halus. Hal tersebut menunjukkan bahwa masyarakat tempo dulu yang konon katanya kebanyakan tidak pernah duduk di bangku sekolah, namun memiliki cara berbahasa yang sangat santun.
\end{quote}

\end{enumerate}

\vspace{0.5cm}

\item Buat/susun daftar pustaka berdasarkan kutipan tersebut dengan benar!

\begin{thebibliography}{9}

\bibitem{}
Kamus Besar Bahasa Indonesia (KBBI). Definisi Peribahasa. Diakses dari \texttt{https://kbbi.kemdikbud.go.id}

\bibitem{}
Liputan6.com. Data jumlah suku bangsa di Indonesia. Diakses dari \texttt{https://www.liputan6.com}

\bibitem{}
Oudri, Nabila. Artikel tentang bahasa daerah yang terancam punah di Indonesia. Diakses dari \texttt{https://itjen.kemdikbud.go.id}

\end{thebibliography}

\vspace{0.5cm}

\item Bagaimana keadaan bahasa daerah kita (Indonesia) berdasarkan isi teks tersebut?
\begin{quote}
    Berdasarkan isi teks, keadaan bahasa daerah di Indonesia cukup memprihatinkan. Dari 718 bahasa daerah yang ada, sebagian besar berada dalam kondisi terancam punah dan kritis. Bahkan, beberapa bahasa daerah telah dinyatakan punah. Hal tersebut menunjukkan bahwa kekayaan bahasa daerah Indonesia mulai mengalami kemunduran dan perlu dilestarikan agar tidak hilang.
\end{quote}
\item Tuliskan peribahasa yang Saudara ketahui selain peribahasa yang terdapat di dalam modul!
\begin{quote}
    \begin{itemize}
        \item \textbf{Kecek berlauk-pauk, makan dengan sambal lada;} membual dengan membangga-banggakan keunggulan (kelebihan) yang dimilikinya, padahal sesungguhnya tiada ia mempunyai. Sombong (takabur).
        \item \textbf{Kecil bernama besar bergelar;} untuk menyebut nama anak muda boleh disebut namanya, namun untuk menyebut nama orang tua hendaklah disebut gelarnya.
        \item \textbf{Kecil-kecil anak, kalau sudah besar onak;} anak ketika kecil menyenangkan hati, namun setelah besar (tindakan atau kelakuannya) menyusahkan hati.
    \end{itemize}
\end{quote}

\end{enumerate}



\begin{thebibliography}{99}
    \bibitem{}
    Yudhistira, Lenin, dan Irawan. 2022. Pembahasan mengenai peribahasa sebagai nilai luhur budaya, halaman 44--45. Diakses dari \url{https://books.google.co.id/books?id=73gBersTfhoC&printsec=copyright&hl=id}

\end{thebibliography}





\end{document}

